%%%% RESUMO
%%
%% Apresentação concisa dos pontos relevantes de um texto, fornecendo uma visão rápida e clara do conteúdo e das conclusões do
%% trabalho.

\begin{resumoutfpr}%% Ambiente resumoutfpr
Com o crescimento da complexidade dos sistemas distribuídos e aplicações modernas, aliada à necessidade de entregas ágeis e confiáveis, surgiu o conceito de \textit{DevOps}, que busca integrar as equipes de desenvolvimento e operações para otimizar a automação, a colaboração e a entrega de \textit{software}. Dentro dessa abordagem, destacam-se as práticas de Integração Contínua (CI) e Entrega Contínua (CD). A CI refere-se ao processo de automação da integração de código de múltiplos desenvolvedores em um repositório compartilhado, garantindo que mudanças frequentes sejam verificadas e testadas continuamente. A CD complementa esse processo, automatizando a entrega do \textit{software} para ambientes de teste e produção, permitindo implantações rápidas e confiáveis. Para que essas práticas sejam eficazes, utilizam-se \textit{pipelines} (ou esteiras de CI/CD), que são conjuntos de processos automatizados responsáveis por compilar, testar e implantar aplicações de maneira sistemática e reprodutível.
Este trabalho apresenta um estudo comparativo entre três das principais ferramentas de CI/CD: GitHub Actions, GitLab CI/CD e Jenkins. A metodologia adotada consiste na implementação de \textit{pipelines} CI/CD em cada plataforma, seguidos de uma avaliação baseada em critérios previamente definidos. Para garantir a confiabilidade dos resultados, os testes serão realizados em um ambiente controlado, onde as variáveis de infraestrutura permanecerão constantes. Com isso, este estudo contribuirá para a escolha da ferramenta de CI/CD mais adequada para diferentes cenários, auxiliando equipes de desenvolvimento na automação de seus fluxos de trabalho.
\end{resumoutfpr}
